\clearpage\thispagestyle{empty}
\pdfbookmark[0]{Quellen zu Über eine Aufgabe der sphärischen Astronomie (1808)}{sphaerisch1808-sources}
\begin{center}\textbf{Quellen zu Über eine Aufgabe der sphärischen Astronomie (1808)}\end{center}
\noindent Diese KI-Übertragung in heutiges Deutsch folgt Carl Friedrich Gauß’ deutschem Original in den \textit{Werken}, Band~VI (Göttingen, 1874), S.~129--140. Der Erstdruck erschien als Beitrag XXXV im Oktoberheft 1808 der von Franz Xaver von Zach herausgegebenen \textit{Monatlichen Correspondenz zur Beförderung der Erd- und Himmels-Kunde}, Band~18, S.~277--293. \href{https://gdz.sub.uni-goettingen.de/id/PPN236007467}{Historischer deutscher Text in den Werken}. \href{https://www.digitale-sammlungen.de/en/view/bsb10538609?page=299}{Erstdruck}.
\par\smallskip\noindent William Chauvenet behandelt Gauß’ praktische Bemerkungen und stellt sein Rechenbeispiel in anderer Form dar, in \href{https://archive.org/details/amanualspherica06chaugoog/page/n291/mode/1up}{\textit{A Manual of Spherical and Practical Astronomy}}, Band~I, \textit{Spherical Astronomy} (Philadelphia und London, 1863), S.~283 und 285.
\par\smallskip\noindent Übertragung und Satz wurden mit KI erstellt.
\clearpage
